\subsection{BBH Subtask Analysis}

\begin{table}[t]
\centering
\small
\setlength{\tabcolsep}{2.5pt}
\begin{tabular}{lcccc}
\toprule
\textbf{Task} & \textbf{Vanilla MoE} & \textbf{Loop Base} & \textbf{\ourmodelnospace}\\
\midrule
hyperbaton            & 31.25 & 50.83 & \textbf{51.25}\\
navigate              & 40.83 & 52.50 & \textbf{57.08}\\
salient trans. err. det. & \textbf{22.50} & 16.67 & 15.00\\
logical deduct. 5 obj. & \textbf{22.50} & 17.50 & 16.25 \\
\bottomrule
\end{tabular}
\caption{Representative BBH subtask scores for Vanilla MoE, Loop Base, and \ourmodelnospace. Here, ``salient trans. err. det.'' and ``logical deduct. 5 obj.'' stand for ``salient translation error detection'' and ``logical deduction five objects'', respectively. Full results across all 27 subtasks are provided in Appendix~\ref{appendix:bbh}.}
\label{tab:bbh_highlight}
\end{table}

The aggregate BBH gain of \ourmodel over the Vanilla MoE conceals a structured heterogeneity across subtasks. This offers finer-grained evidence on where iterative depth contributes most. Table~\ref{tab:bbh_highlight} reports the two subtasks with the largest gains and the two with the largest regressions.

The two largest improvements, hyperbaton ($+20.00$) and navigate ($+16.25$), both require answers to be constructed through a sequence of dependent reasoning steps in which each step builds on context accumulated from prior iterations. The Loop Base alone already accounts for the bulk of these gains, capturing nearly the entire improvement of hyperbaton and roughly three-quarters of the improvement of navigate. This pattern indicates that iterative depth is the dominant mechanism, while IterAdaLN and capacity balancing provide a smaller secondary contribution on tasks whose reasoning chains benefit from per-token modulation. The finding offers subtask-level corroboration of the main result that loop iterations supply the additional computation required for complex multi-step reasoning. 

The two largest regressions, salient translation error detection and five-object logical deduction, exhibit a monotonic decline that is already present in Loop Base. We attribute this to reduced per-step activation breadth, since matching the Vanilla MoE in compute forces LoopMoE to activate fewer physical parameters per iteration.

Across the full 27-subtask distribution, gains and regressions align consistently with this principle. \ourmodel improves on tasks where reasoning is compositional and benefits from iterative refinement and regresses on tasks whose correctness depends on broad, one-shot information access within each step.




\subsection{Iterative Routing Dynamics}
\label{sec:routing_analysis}

Beyond aggregate task performance, we next examine the iterative structure of a looped MoE block, which raises a natural question about expert allocation, namely, whether successive iterations activate similar or dissimilar expert subsets. Activating largely disjoint subsets would let each iteration perform a functionally distinct computation and effectively expand the active parameter budget, whereas converging to the same subset would form a stable expert coalition that applies a consistent transformation under repeated application. To probe which regime the loop operates in, we measure two pairwise similarities between iterations. Cosine similarity of router inputs tracks continuous representation drift, and Jaccard similarity of selected expert sets tracks discrete routing change.

\begin{figure}[t]
    \centering
    \includegraphics[width=0.48\columnwidth]{Images/loop_cosine_summary.png}
    \includegraphics[width=0.48\columnwidth]{Images/loop_jaccard_summary.png}
    \caption{Cross-layer pairwise similarity between loop iterations on the BBH dataset. Router-input cosine (left) and expert-activation Jaccard (right). Per-layer matrices are reported in Appendix~\ref{app:per_layer_routing}.}
    \label{fig:loop_dynamics}
\end{figure}

The heatmaps in Figure~\ref{fig:loop_dynamics} reveal a three-phase structure that interpolates between the two extremes at different points in the loop. In the entry phase, iter~0 stands apart from all subsequent iterations and is especially distant in cosine similarity. This is the only iteration that reads directly from the non-loop prefix rather than from the previous block iteration, and it behaves as an adapter that projects the prefix distribution into the recurrent state. In the recurrent core, iters~1 and~2 form the tightest cluster and reach the highest pairwise Jaccard similarity, approaching the fixed-coalition regime. Once the prefix has been absorbed, the loop settles into a stationary internal recurrence carried by a consistent expert subset~\cite{blayney2026mechanistic}. In the exit phase, the terminal iteration exhibits a notable decoupling between the two metrics. Representations remain close to those of the previous iteration, yet the router selects a substantially different expert subset, which suggests that the final pass reformats the consolidated state for the downstream non-loop layers rather than continuing to refine it. A per-layer decomposition that shows how the two loop layers contribute asymmetrically to this structure is provided in Appendix~\ref{app:per_layer_routing}.
